\documentclass[11pt]{article}
\usepackage[a4paper,margin=26mm]{geometry}
\usepackage{amsmath,amssymb}
\usepackage[hidelinks]{hyperref}
\setlength{\emergencystretch}{3em}
\title{Conjecture 00000001149: an impossible norm for a Gaussian coefficient}
\author{lizaixi01, with OpenAI Codex assistance}
\date{22 October 2026}
\begin{document}
\maketitle
\section{The necessary condition already fails}
The conjecture asks for a $p$-adically convergent series whose $m$th
term is a $q$-binomial coefficient with absolute value $p^{-m/2}$,
for coprime minimal-model parameters $p,q$.
Take $p=2$ and $q=5$. These satisfy $2\leq p<q$ and
$\gcd(p,q)=1$. Regardless of which fusion coefficient is to be
represented, the term with $m=1$ would have to satisfy
\[
 |a_1|_2=2^{-1/2}.
\]
This is impossible for any $a_1\in\mathbf Q_2$, hence in particular
for a Gaussian coefficient at the integer parameter $q=5$.
The argument applies whether the series is numbered from zero or one,
since both conventions include the term indexed by one.

Indeed, every nonzero $x\in\mathbf Q_2$ has an integer valuation
$v_2(x)$ and $|x|_2=2^{-v_2(x)}$. If this were $2^{-1/2}$,
injectivity of the real exponential with base $2$ would give
$2v_2(x)=1$, a contradiction. Zero also cannot have that norm.
Thus no sequence meeting the asserted termwise condition exists;
questions about its convergence and its sum do not arise.

\section{The coefficients are actual rational expressions}
We use the standard Gaussian coefficient, not a free variable
assigned a convenient type. In the notation of
\href{https://dlmf.nist.gov/17.2.E27}{NIST DLMF, equation 17.2.27},
\[
 (q;q)_n=\prod_{j=0}^{n-1}(1-q^{j+1}),\qquad
 \genfrac{[}{]}{0pt}{}{n}{k}_q=\frac{(q;q)_n}{(q;q)_k(q;q)_{n-k}}
 \quad(0\leq k\leq n).
\]
At $q=5$ all factors in the denominator are nonzero rational
numbers, so every displayed coefficient is rational. We set the
coefficient to zero for $k>n$, which cannot help achieve the required
nonzero norm either. Integrality of Gaussian coefficients is a
stronger classical fact, but is unnecessary for this obstruction.

The Lean source defines these finite products and quotients over
$\mathbf Q$, proves their denominators nonzero, and maps the result
into the actual field \texttt{Padic 2}. It also considers the canonical
image in \texttt{PadicAlgCl 2}; the extended norm agrees with the
$\mathbf Q_2$ norm. Therefore passing to the algebraic closure does
not alter the norm of any of these coefficients.

\newpage
\section{Formal statement and scope}
\texttt{no\_half\_integral\_norm} proves the obstruction for every
element of $\mathbf Q_2$. The Gaussian-coefficient definitions then
give \texttt{gaussianPadic\_ne\_half\_norm} and its algebraic-closure
version. Finally, \texttt{no\_requested\_gaussian\_series} quantifies
over arbitrary natural-number upper and lower index sequences
$n(m),k(m)$ and rules out the required norm law for all $m\geq1$.
The theorem \texttt{counterexample} combines this with the explicit
legal parameter check for $(2,5)$.

This disproves a necessary clause of the question as printed.
Here $q$ is the integer minimal-model parameter appearing in the
same statement, and the coefficients are the standard, unscaled
Gaussian coefficients. In conformal-field-theory literature the
letter $q$ may instead denote a separate nome. That alternative
parameter convention is not specified in this question and is
not asserted to be covered by this argument.
It neither defines nor evaluates the Virasoro fusion coefficients,
and makes no claim about a modified norm law, a different meaning
of a generalized $q$-binomial coefficient, or arbitrary algebraic
numbers outside the rational-coefficient image. The obstruction
does not assert that $2^{-1/2}$ is absent from the norm-value set of
the entire algebraic closure.

\section{Verification and provenance}
The accompanying project pins Lean 4.33.0 and Mathlib commit
\[
\texttt{db584cd6d46c92f209a44c0f1c829460d327499d}.
\]
The source-bound build, transitive axiom audit and official fresh
kernel replay are recorded in the package evidence. No
\texttt{sorry}, \texttt{native\_decide}, or extra axiom is used.
The source statement is preserved at its upstream commit in
\texttt{SOURCE.md}. OpenAI Codex assisted the proof and report;
a separate supervising Codex agent reviewed the formal chain.
This is AI review, not human certification. Organizer acceptance
is pending.
\end{document}
